% ============================================================================
% Section 2 -- background and system model.
%
% STATUS: first full draft, 2026-09-05.
%
% This section carries the paper's largest reviewer risk (plan section 6.1).
% Riders are admitted with ZERO main-pool blocks and their prompt declared
% already computed. That is only legitimate on a decode node whose KV arrives
% from a prefill node. Stated here it is a scoping assumption about the
% architecture the field already runs; discovered by a reviewer in the
% evaluation it reads as concealment.
% ============================================================================

\section{Background and system model}
\label{sec:background}

\subsection{Capacity, not compute, caps the batch}

A decode step is memory-bound by a wide margin. It reads every model weight
once and every resident KV byte once, and performs arithmetic whose intensity
equals the GQA group size --- 4 to 8 on current models, against a balance point
near 116 FLOP/byte on an H100. The consequence is that a decode engine runs at
single-digit arithmetic utilization, and what stops it running a larger batch
is not the arithmetic but the bytes: KV cache exhausts the device long before
the SMs saturate. A recent trace-driven study of production serving reaches the
same conclusion from the workload side, finding that ``VRAM capacity, not token
budget'' is what binds~\cite{bottlenecks}.

It is worth being precise about which regime this is, because the opposite
claim is also true somewhere. At small batch the step is dominated by the
weight read and adding requests is nearly free --- decode is
\emph{weight-bandwidth-bound}. That regime is real, but it is precisely the
regime in which capacity is not a problem, so it is not the one a capacity
feature should be designed for. At the operating point that matters, where the
pool is full and the scheduler is turning requests away, the step is roughly
79\% KV read and 21\% weight read. Everything in this paper concerns that
second regime, and \S\ref{sec:model} makes the boundary between them explicit.

\subsection{Host memory as a tier for actively decoding requests}

Offloading KV to host memory is an old idea that a new interconnect has
changed. Over PCIe~4 at roughly 25\,GB/s, streaming a request's KV to the GPU
every decode step costs more than the step itself, so host memory could only
ever be a \emph{swap} target: somewhere to park a request that is not currently
running. NVLink-C2C moves 450\,GB/s in each direction,\footnote{The commonly
quoted 900\,GB/s is bidirectional. A KV stream is unidirectional
host-to-device, so 450\,GB/s is the figure that bounds it, and we measure
rather than assume even that.} and at that rate the bytes a decoding request
needs can cross the link inside the time the GPU spends reading its own
weights.

That changes what host memory \emph{is}. It becomes a tier that an
\emph{actively decoding} request can occupy, which is a state that did not
previously exist and which no existing KV policy governs --- continuous
batching, prefix reuse, eviction and hierarchical caching all decide the
residency of blocks belonging to requests that are not currently running.

\subsection{System model: a disaggregated decode node}
\label{sec:background:pd}

We target a \textbf{decode node in a prefill/decode-disaggregated
deployment}~\cite{distserve,splitwise,mooncake}. Prefill runs on separate
hardware; a request arrives at the decode node with its KV cache already
computed, delivered over the cluster transport. This is the production default
rather than a convenience we adopt: the separation exists because prefill and
decode have opposite resource profiles, and every major serving stack now
supports it.

The assumption does real work here, and we state exactly what it buys. On a
decode node, a host-tier request's KV can be delivered \emph{into host memory}
and never occupy an HBM block at all --- not even transiently. Riders are
therefore \textbf{additive}: they are served in addition to a full resident
cohort rather than displacing part of it.

In a combined prefill+decode engine this is false, and the difference is not
subtle. A request must be prefilled somewhere, so its KV lands in HBM first and
is only reclaimed afterwards; riders and residents then draw on the same
headroom, and we measured the result --- the served total creeps toward the
all-HBM capacity from below and never passes it, at any setting of the
resident-cohort cap. A rider whose prefill stays resident frees nothing while
still costing interconnect bandwidth every step, which is strictly worse than
leaving it in HBM. We report that negative result in
\S\ref{sec:eval:capacity}, because it is what motivates the mechanism and
because it bounds where the technique applies.

The assumption's cost is that KV must reach host DDR. We do not wave this
away: \S\ref{sec:impl:stash} implements the handover and \S\ref{sec:eval}
reports what it costs, so the reader can price the assumption rather than
accept it.

% NOTE: a subsection defining "lossless" and describing how we verify it was
% removed 2026-09-06. Exact placement is a property of moving bytes rather than
% approximating them, and no comparable systems paper argues for it -- doing so
% invites the doubt it was meant to settle. The bf16 near-tie methodology, which
% is a real gotcha, now lives as one paragraph in the evaluation setup where it
% belongs.
